% Part: normal-modal-logic
% Chapter: axioms-systems
% Section: normal-logics

\documentclass[../../../include/open-logic-section]{subfiles}

\begin{document}

\olfileid{nml}{prf}{nor}

\olsection{સામાન્ય મોડલ તર્કશાસ્ત્રો}

મોડલ !!{formula}sના દરેક ગણને સ્વયંસિદ્ધિઓના કોઈ ગણમાંથી નિષ્પન્ન થતા
!!{formula}sના ગણ તરીકે સહેલાઈથી લક્ષણિત કરી શકાતો નથી. આપણે મોડલ તર્કો
સુવ્યવસ્થિત હોય એવું ઇચ્છીએ છીએ. સૌપ્રથમ, પરંપરાગત વિધાનાત્મક તર્કમાં
જે કંઈ આપણે !!{derive} કરી શકીએ તે હજી પણ !!{derivable} હોવું જોઈએ;
અલબત્ત, !!{formula}sમાં હવે \iftag{prvBox}{$\Box$\iftag{prvDiamond}{
    અને~}{}}{}\iftag{prvDiamond}{$\Diamond$}{} પણ હોઈ શકે છે. આ માટે
અમે માગીએ છીએ કે મોડલ તર્કમાં પુનરુક્તિઓના બધા દૃષ્ટાંતો હોય અને તે
મોડસ પોનેન્સ હેઠળ બંધ હોય.

\begin{defn}
  \emph{મોડલ તર્કશાસ્ત્ર} એ મોડલ !!{formula}sનો એવો ગણ~$\Sigma$ છે જે
  \begin{enumerate}
  \item બધી પુનરુક્તિઓ ધરાવે છે, અને
  \item પ્રતિસ્થાપન હેઠળ બંધ છે, એટલે કે જો $!A \in \Sigma$ હોય અને
    $!D_1$, \dots, $!D_n$ !!{formula}s હોય, તો
    \[
    \SSubst{!A}{\subst{!D_1}{p_1}, \dots, \subst{!D_n}{p_n}} \in \Sigma,
    \]
    \item \emph{મોડસ પોનેન્સ} હેઠળ બંધ છે, એટલે કે જો $!A$ અને $!A
      \lif !B \in \Sigma$ હોય, તો $!B \in \Sigma$ હોય છે.
  \end{enumerate}
\end{defn}

મોડલ તર્કો માટે સંબંધાત્મક અર્થવિજ્ઞાન વાપરવું હોય, તો બધાં મોડલ
નિદર્શનોમાં પ્રામાણ્ય હોય એવાં બધાં !!{formula}s પણ તેમાં સામેલ હોય તેવી
માગ રાખવી પડે. \Ax{K}\iftag{prvDiamond}{ અને~\Dual{}}{}ના બધા દૃષ્ટાંતો
!!{derivable} હોય, અને !!a{formula}~$!A$ !!{derivable} હોય ત્યારે
$\Box !A$ પણ નિષ્પન્ન હોય, તો આ માગ પૂરી થાય છે. આ શરતો પૂરી કરતો
મોડલ તર્ક \emph{સામાન્ય} કહેવાય છે. (નિશ્ચિતપણે, અસામાન્ય મોડલ તર્કો પણ
છે, પરંતુ સામાન્ય સંબંધાત્મક નિદર્શનો તેમને માટે પર્યાપ્ત નથી.)

\begin{defn}
  મોડલ તર્ક $\Sigma$ \emph{સામાન્ય} છે, જો તે આ સૂત્રો ધરાવે:
  \begin{align*}
    \tag{\Ax{K}} & \Box(p \lif q) \lif (\Box p \lif \Box q),
    \iftag{prvDiamond}{\\
      \tag{\Dual} & \Diamond p \liff \lnot\Box\lnot p}{}
  \end{align*}
  અને તે \emph{આવશ્યકીકરણ} હેઠળ બંધ હોય, એટલે કે $!A \in
  \Sigma$ હોય, તો $\Box !A \in \Sigma$ પણ હોય.
\end{defn}

નોંધો કે પુનરુક્તિમૂલક પ્રેરણ ``સૂક્ષ્મ'' અર્થમાં કોઈ \emph{જગતમાં
સત્યતા} જાળવવા જેટલું પૂરતું છે, જ્યારે નિયમ \Nec{} માત્ર \emph{નિદર્શનમાં
સત્યતા} જાળવે છે (અને તેથી ફ્રેમમાં કે ફ્રેમોના વર્ગમાં પ્રામાણ્યતા પણ
જાળવે છે).

\begin{prop}\ollabel{prop:rk}
  દરેક સામાન્ય મોડલ તર્કશાસ્ત્ર નિયમ \RK હેઠળ બંધ છે:
  \begin{prooftree}
    \AxiomC{$!A_1 \lif (!A_2 \lif \cdots (!A_{n-1} \lif !A_n)\cdots)$}
    \RightLabel{\RK}
    \UnaryInfC{$\Box!A_1 \lif (\Box!A_2 \lif \cdots (\Box!A_{n-1}
      \lif \Box!A_n)\cdots).$}
  \end{prooftree}
\end{prop}

\begin{proof}
  $n$ પર આગમનથી: $n = 1$ હોય, તો આ નિયમ માત્ર \Nec છે, અને દરેક
  સામાન્ય મોડલ તર્ક \Nec હેઠળ બંધ છે.

  હવે માનો કે પરિણામ $n-1$ માટે સાચું છે; આપણે તે~$n$ માટે બતાવીએ.

  માનો કે
  \begin{align*}
  & !A_1 \lif (!A_2 \lif \cdots (!A_{n-1} \lif !A_n)\cdots) \in \Sigma
  \intertext{આગમન-પરિકલ્પના મુજબ આપણને મળે છે}
  & \Box!A_1 \lif (\Box!A_2 \lif \cdots \Box(!A_{n-1} \lif !A_n)\cdots)
  \in \Sigma
  \intertext{$\Sigma$ સામાન્ય મોડલ તર્ક હોવાથી તે~$\Ax{K}$ના બધા દૃષ્ટાંતો,
    ખાસ કરીને આ દૃષ્ટાંત, ધરાવે છે}
  & \Box(!A_{n-1} \lif !A_n) \lif (\Box!A_{n-1} \lif \Box!A_n) \in \Sigma
  \intertext{મોડસ પોનેન્સ અને પુનરુક્તિઓના યોગ્ય દૃષ્ટાંતો વાપરવાથી મળે છે}
  & \Box!A_1 \lif (\Box!A_2 \lif \cdots (\Box!A_{n-1}
  \lif \Box!A_n)\cdots) \in \Sigma. 
  \end{align*}
\end{proof}

\begin{prop}\ollabel{prop:notDiamondBot}
  દરેક સામાન્ય મોડલ તર્ક $\Sigma$માં~$\lnot\Diamond\lfalse$ હોય છે.
\end{prop}

\begin{prob}
  \olref[nml][prf][nor]{prop:notDiamondBot} સાબિત કરો.
\end{prob}

\begin{prop}
  માનો કે $!A_1$, \dots, $!A_n$ !!{formula}s છે. ત્યારે $!A_1$,
  \dots, $!A_n$ના બધા દૃષ્ટાંતો ધરાવતો સૌથી નાનો સામાન્ય મોડલ તર્ક
  $\Sigma$ છે.
\end{prop}\sourcecorrection{OLMD-023}{મૂળ પ્રસ્તાવમાં માત્ર ``મોડલ તર્ક'' લખાયું છે, પરંતુ તેની સાબિતી સામાન્ય મોડલ તર્કોના છેદથી સૌથી નાનો સામાન્ય મોડલ તર્ક રચે છે; આ વિશેષણ અહીં સ્પષ્ટ કર્યું છે.}

\begin{proof}
  $!A_1$, \dots, $!A_n$ આપેલાં હોય ત્યારે $\Sigma$ને એવાં બધાં સામાન્ય
  મોડલ તર્કોના છેદ તરીકે વ્યાખ્યાયિત કરો જે $!A_1$, \dots, $!A_n$ના
  બધા દૃષ્ટાંતો ધરાવે છે. આવો છેદ ખાલી નથી, કારણ કે બધાં !!{formula}sનો
  ગણ~$\Frm[L]$ પોતે એવો મોડલ તર્ક છે.
\end{proof}

\begin{defn}
$!A_1$, \dots, $!A_n$ ધરાવતો સૌથી નાનો સામાન્ય મોડલ તર્ક
\emph{મોડલ તંત્ર} કહેવાય છે અને તેને $\Log{K} !A_1 \dots
!A_n$ વડે દર્શાવાય છે. સૌથી નાનો સામાન્ય મોડલ તર્ક~\Log{K} વડે
દર્શાવાય છે.
\end{defn}

\end{document}
